% ============================================================
%  CAPITOLO 4 – Una Visione Integrata del Patrimonio
%               Previdenziale
% ============================================================
\chapter{Una Visione Integrata del Patrimonio Previdenziale}
\label{cap:patrimonio_integrato}

Il capitolo precedente ha mostrato che la previdenza complementare è una
necessità strutturale per chiunque voglia evitare la povertà relativa in età
anziana. Ma affidarsi esclusivamente al fondo pensione sarebbe un errore
simmetrico a quello di affidarsi esclusivamente al primo pilastro pubblico:
entrambe le scelte riducono il portafoglio previdenziale a una singola
scommessa. Il presente capitolo allarga la prospettiva e propone un
framework integrato in cui la pensione non è il prodotto di un'unica scelta,
ma di un insieme coerente di decisioni di vita, carriera e risparmio che si
compongono lungo quarant'anni.

La tesi centrale è la seguente: il \emph{patrimonio previdenziale} di un
lavoratore della Generazione Z non è un numero fisso determinato dalla
formula NDC, ma una funzione di almeno tre classi di variabili: il capitale
umano accumulato (che determina il sentiero salariale e quindi il montante
contributivo), la diversificazione del risparmio privato tra i tre pilastri, e
il livello di consapevolezza e di educazione finanziaria che permette di
governare le prime due con razionalità. Queste tre dimensioni sono
analizzate nelle tre sezioni che seguono.

% ============================================================
\section{Il capitale umano come primo asset previdenziale}
\label{sec:capitale_umano}
% ============================================================

Gary Becker, nella sua opera seminale \emph{Human Capital} del 1964,
introdusse una rivoluzione concettuale che l'economia previdenziale ha
recepito in ritardo: l'investimento in istruzione e formazione non è un costo
di consumo ma l'acquisizione di un asset produttivo, il cui rendimento si
materializza lungo l'intera vita lavorativa sotto forma di salari più alti,
carriere più stabili e minore esposizione alla disoccupazione
frizionale~\cite{Becker1964}. Nel contesto del sistema NDC italiano, questo
principio acquista una valenza previdenziale diretta: il montante contributivo
è proporzionale ai salari versati, dunque ogni euro investito in formazione
che produce un aumento permanente del salario genera un effetto moltiplicativo
sul montante finale che si protrae per decenni.

\subsection{Istruzione, competenze e montante contributivo:
            il capitale umano nell'equazione NDC}
\label{subsec:istruzione_ndc}

L'equazione del montante contributivo sviluppata nel Capitolo~1 può essere
riscritta evidenziando esplicitamente la dipendenza dal profilo salariale
individuale. Se il salario nell'anno di lavoro \(j\) cresce al tasso
\(m > 0\) a partire da un salario iniziale \(S_0\), si ha
\(S_j = S_0 \cdot (1+m)^{j-1}\), e il montante contributivo al
pensionamento dopo \(L\) anni diventa:
%
\begin{equation}
  \mathrm{MC}_{\mathrm{NDC}}(m)
  = c \cdot S_0
    \sum_{j=1}^{L} (1+m)^{j-1}(1+g)^{L-j}
  \label{eq:mc_m}
\end{equation}
%
Dalla formula~\eqref{eq:mc_m} emerge il ruolo del tasso \(m\) di crescita
salariale come \emph{secondo moltiplicatore} del montante, accanto al
tasso di capitalizzazione nozionale \(g\)~\cite{Bosi2018}. La letteratura
sul capitale umano, a partire dalle equazioni salariali di Mincer (1974),
documenta sistematicamente che i lavoratori con istruzione superiore
presentano tassi di crescita salariale \(m\) strutturalmente più elevati di
quelli con istruzione secondaria: secondo le stime OCSE per l'Italia,
il differenziale di crescita salariale tra un laureato magistrale e un
diplomato supera i 2--3 punti percentuali annui nel lungo periodo~\cite{Mincer1974, OECD2023}.

\paragraph{Esempio numerico comparato.}
Per quantificare l'impatto previdenziale di questo differenziale, si
confrontano due profili a parità di tutti gli altri parametri
(\(c = 0{,}33\), \(S_0 = 25.000\) euro, \(g = 2{,}7\%\), \(L = 40\) anni,
\(e(67) = 20\) anni):

\begin{itemize}
  \item \textbf{Profilo A} -- lavoratore con basso investimento in capitale
    umano: \(m_A = 1{,}5\%\) annuo.
  \item \textbf{Profilo B} -- lavoratore con alta specializzazione
    (laurea magistrale, competenze tech-economiche): \(m_B = 4{,}0\%\) annuo.
\end{itemize}

Utilizzando la formula~\eqref{eq:mc_m} con l'approssimazione della
somma geometrica:
%
\begin{align}
  \mathrm{MC}(m_A) &= c \cdot S_0 \cdot
    \frac{(1+m_A)^L - (1+g)^L}{m_A - g}
    \approx 0{,}33 \times 25.000 \times 98{,}4
    \approx \text{\euro}\,811.800
    \label{eq:mc_A}\\[4pt]
  \mathrm{MC}(m_B) &= c \cdot S_0 \cdot
    \frac{(1+m_B)^L - (1+g)^L}{m_B - g}
    \approx 0{,}33 \times 25.000 \times 148{,}3
    \approx \text{\euro}\,1.223.475
    \label{eq:mc_B}
\end{align}

Applicando il coefficiente di trasformazione corrente a 67 anni
(\(\gamma = 5{,}723\%\)), la pensione annua lorda del Profilo B supera quella
del Profilo A di circa \(\text{\euro}\,23.700\) l'anno --
un divario di oltre \(\text{\euro}\,475.000\) nell'arco dell'intera fase di
quiescenza, ipotizzando 20 anni di percezione.

\begin{table}[htbp]
\centering
\caption{Impatto del tasso di crescita salariale \(m\) sul montante
         e sulla pensione NDC\\
         (\(c=0{,}33\), \(S_0=\text{\euro}\,25.000\), \(g=2{,}7\%\),
          \(L=40\), \(e(67)=20\) anni)}
\label{tab:montante_m}
\begin{tabular}{lrrrr}
\toprule
Profilo & \(m\) & Montante lordo &
  Pensione annua lorda & Tasso sostituz.\ approx.\\
\midrule
A -- basso cap.\ umano & \(1{,}5\%\) &
  \euro\,811.800 & \euro\,46.456 & $\approx 58\%$ \\
B -- alto cap.\ umano  & \(4{,}0\%\) &
  \euro\,1.223.475 & \euro\,70.011 & $\approx 70\%$ \\
\midrule
Differenziale B\,$-$\,A & $+2{,}5\,pp$ &
  +\euro\,411.675 & +\euro\,23.555 & $+12\,pp$ \\
\bottomrule
\end{tabular}
\smallskip\\
\footnotesize Elaborazione propria su formula~\eqref{eq:mc_m}.
Parametro \(g\) = scenario base DFP 2025~\cite{MEFRGS2025}.
\end{table}

La Tabella~\ref{tab:montante_m} dimostra che la scelta di investire in
istruzione di alta qualità non è, per il lavoratore NDC puro, soltanto una
strategia di carriera: è letteralmente un'operazione di pianificazione
previdenziale che sposta il tasso di sostituzione atteso dal 58\% al 70\%,
soglia di mantenimento del tenore di vita identificata dalla letteratura~\cite{Bosi2018}.
Non è narrativa motivazionale: è l'equazione NDC applicata alla scelta di
investimento in istruzione~\cite{Becker1964, Mincer1974}.

\subsection{Il Lifelong Learning come strategia previdenziale:
            occupabilità continua uguale continuità contributiva
            nell'era dell'automazione}
\label{subsec:lifelong_learning}

L'argomento del capitale umano come asset previdenziale potrebbe essere
interpretato come una raccomandazione statica: investi nell'istruzione
iniziale e il resto viene da sé. Ma nell'economia del XXI secolo,
caratterizzata da cicli tecnologici sempre più rapidi, questa lettura è
pericolosamente obsoleta. Acemoglu e Restrepo (2018) hanno documentato
che l'automazione non distrugge uniformemente l'occupazione, ma la
\emph{ridistribuisce}: crea nuovi compiti ad alta complementarità con la
tecnologia e obsoleta quelli routinari e codificabili~\cite{AcemogluRestrepo2018}.
La velocità di questa ridistribuzione si è accelerata nell'era dell'intelligenza
artificiale generativa: secondo le stime del World Economic Forum (2025),
il 40\% delle competenze attualmente richieste dal mercato del lavoro sarà
parzialmente o completamente obsoleto entro il 2030~\cite{WEF2025}.

La conseguenza per la previdenza è diretta. Nel sistema NDC, la continuità
contributiva è la condizione necessaria per l'adeguatezza della prestazione
finale. Un lavoratore che perde la propria occupabilità a 45 anni -- perché
le sue competenze sono state sostituite dall'automazione -- e impiega due
anni a riqualificarsi, subisce una doppia perdita: la perdita contributiva
diretta (due anni di contributi mancanti al montante) e la perdita salariale
indiretta (rientro nel mercato del lavoro a un livello salariale più basso).
Il \emph{lifelong learning} -- la formazione continua lungo tutto l'arco della
vita lavorativa -- non è quindi un lusso per i più ambiziosi: è la condizione
necessaria per mantenere l'occupabilità e, con essa, la continuità
contributiva~\cite{Bosi2018, Raitano2022}.

Formalmente, se definiamo \(\tau\) come il numero di anni di gap contributivo
dovuto a disoccupazione tecnologica non anticipata, il costo previdenziale
del mancato lifelong learning è:
%
\begin{equation}
  \Delta \mathrm{MC}(\tau) =
  c \sum_{j=j^*}^{j^*+\tau} S_j (1+g)^{L-j}
  \label{eq:costo_gap}
\end{equation}
%
dove \(j^*\) è il primo anno di interruzione. Per un lavoratore con
\(S_{j^*} = 35.000\) euro, \(\tau = 2\), \(g = 2{,}7\%\) e
\(L - j^* = 20\) anni residui, la perdita di montante
stimata dalla~\eqref{eq:costo_gap} supera i \euro\,60.000. Questa cifra
non tiene conto della perdita salariale successiva: il costo totale della
mancata formazione è sistematicamente sottostimato dai lavoratori, perché
si materializza quarant'anni dopo e non produce nessun segnale di feedback
immediato~\cite{ThalerSunstein2008}.

\subsection{Le carriere discontinue e la ``trappola contributiva'':
            come i gap da formazione, maternità o disoccupazione
            erodono il montante}
\label{subsec:trappola_contributiva}

Il capitolo precedente ha già documentato l'impatto delle carriere
discontinue sulla pensione NDC. Vale la pena approfondire qui la
\emph{geometria} della trappola contributiva, che rende i gap all'inizio
della carriera strutturalmente più costosi di quelli alla fine. Dalla
formula~\eqref{eq:mc_m}, il peso di un contributo versato nell'anno \(j\)
sul montante finale è proporzionale al fattore di capitalizzazione
\((1+m)^{j-1}(1+g)^{L-j}\). Per i primi anni di carriera (\(j\) piccolo),
il secondo termine \((1+g)^{L-j}\) è massimo: ogni contributo mancato
in giovinezza è capitalizzato per l'intera durata residua della carriera.

La Tabella~\ref{tab:trappola} quantifica questa asimmetria per Marco,
il lavoratore-tipo della presente tesi, che entra nel mercato del lavoro
a 29 anni nel 2022 e si aspetta di andare in pensione a 67 anni nel 2060
(\(L = 38\), \(g = 2{,}7\%\), \(c = 0{,}27\), \(S_0 = 25.000\) euro,
\(m = 2\%\)).

\begin{table}[htbp]
\centering
\caption{Costo previdenziale di un anno di gap contributivo
         in funzione dell'anzianità di carriera -- caso Marco
         (\(S_0=\text{\euro}\,25.000\), \(m=2\%\), \(g=2{,}7\%\), \(L=38\))}
\label{tab:trappola}
\begin{tabular}{lcccc}
\toprule
Momento del gap & Anno \(j\) &
  Salario \(S_j\) & Fattore cap.\ \((1+g)^{L-j}\) &
  Perdita di montante \\
\midrule
Inizio carriera  & 1  & \euro\,25.000 & $(1{,}027)^{37} = 2{,}68$ &
  $-$\euro\,18.090 \\
Metà carriera    & 19 & \euro\,35.716 & $(1{,}027)^{19} = 1{,}66$ &
  $-$\euro\,15.836 \\
Fine carriera    & 37 & \euro\,49.975 & $(1{,}027)^{1}  = 1{,}03$ &
  $-$\euro\,13.793 \\
\bottomrule
\end{tabular}
\smallskip\\
\footnotesize \(c = 0{,}27\) (aliquota media ibrida). Perdita = contributo
mancato \(\times\) fattore di capitalizzazione.
Elaborazione propria~\cite{Bosi2018, MEFRGS2025}.
\end{table}

I risultati della Tabella~\ref{tab:trappola} evidenziano che un anno di
gap contributivo all'inizio della carriera vale il 31\% in più di un anno
di gap alla fine. Per le lavoratrici che interrompono per maternità
tipicamente intorno ai 30--35 anni, il costo previdenziale di ogni anno
di interruzione si situa nella fascia più costosa. Il CERP ha documentato
come le interruzioni da cura non retribuita contribuiscano a spiegare fino
al 40\% del gender pension gap osservato per le coorti in regime contributivo
puro~\cite{Raitano2022, ForneroMonticone2011}. La ``trappola contributiva''
non è, in ultima analisi, il prodotto di scelte individuali inefficienti: è
il prodotto di strutture del mercato del lavoro e di politiche di welfare che
non riconoscono il valore economico del lavoro di cura. Ignorarla significa
accettare che il sistema NDC, neutrale in apparenza, amplifichi le
disuguaglianze di genere e di classe già presenti nel mercato del lavoro.

% ============================================================
\section{Diversificazione del patrimonio a lungo termine}
\label{sec:diversificazione}
% ============================================================

La previdenza pubblica (primo pilastro) e la previdenza complementare
(secondo pilastro) non esauriscono l'insieme degli strumenti con cui un
individuo può costruire la propria sicurezza economica nella vecchiaia. Esiste
un terzo livello -- il risparmio privato individuale non vincolato -- che
la letteratura chiama convenzionalmente \emph{terzo pilastro informale}
o \emph{pilastro zero}, a seconda del framework di riferimento~\cite{WorldBank2005}.
Questo terzo livello non ha i vantaggi fiscali del fondo pensione, ma offre
flessibilità totale di accesso e una diversificazione dei rischi che i
primi due pilastri non sono in grado di garantire da soli.

\subsection{Il terzo pilastro informale:
            ETF, piani di accumulo (PAC) e risparmio privato
            come complemento al fondo pensione}
\label{subsec:terzo_pilastro}

Il risparmio privato a lungo termine si è trasformato radicalmente
nell'ultimo decennio grazie alla disponibilità di strumenti a basso costo
come gli Exchange Traded Funds (ETF) su indici globali e i Piani di
Accumulo del Capitale (PAC) automatizzati. Un ETF su indice MSCI World --
che replica la performance di circa 1.500 aziende in 23 paesi sviluppati --
ha un costo di gestione (Total Expense Ratio, TER) tipicamente compreso
tra lo 0,10\% e lo 0,20\% annuo, rispetto all'1--2\% dei Piani Individuali
Pensionistici (PIP) assicurativi~\cite{COVIP2024}.

La differenza di costo si traduce, su orizzonti di quarant'anni, in
una differenza di montante finale di portata macroeconomica individuale.
Per un piano di accumulo mensile di \euro\,200, con rendimento lordo del 6\%
annuo e durata 40 anni:
%
\begin{align}
  M_{\text{ETF}}  &= 200 \cdot \frac{(1+r_e)^{480} - 1}{r_e}
    \quad r_e = \frac{0{,}06 - 0{,}002}{12} \approx 0{,}00483
    \quad\Longrightarrow\quad M_{\text{ETF}} \approx \text{\euro}\,388.000
    \label{eq:pac_etf}\\[4pt]
  M_{\text{PIP}}  &= 200 \cdot \frac{(1+r_p)^{480} - 1}{r_p}
    \quad r_p = \frac{0{,}06 - 0{,}015}{12} \approx 0{,}00375
    \quad\Longrightarrow\quad M_{\text{PIP}} \approx \text{\euro}\,283.000
    \label{eq:pac_pip}
\end{align}

Il differenziale è di circa \euro\,105.000, ovvero il 37\% del montante
del PAC-ETF, generato esclusivamente dalla differenza di costo e non da
scelte di rischio~\cite{COVIP2024}. Questo calcolo non tiene conto della
tassazione sui capital gain del PAC-ETF (26\% sul rendimento, anziché
il 12,5\% per i titoli di Stato), che riduce leggermente il vantaggio
netto ma non ne muta la direzione per rendimenti azionari di lungo periodo.

È tuttavia necessario sottolineare la differenza strutturale rispetto al
fondo pensione. Il PAC-ETF non gode delle agevolazioni fiscali del secondo
pilastro (deducibilità dei contributi, tassazione agevolata al 9--15\% in
fase di erogazione, contributo del datore di lavoro): il vantaggio del fondo
pensione per chi ha accesso a queste componenti rimane superiore. Il terzo
pilastro informale è un complemento -- non un sostituto -- del secondo
pilastro istituzionale~\cite{Bosi2018, WorldBank2005}. La sua utilità è
massima per lavoratori autonomi senza accesso ai fondi negoziali, per chi
ha già saturato il limite di deducibilità del fondo pensione
(\euro\,5.164 annui), e per chi vuole costruire una riserva di liquidità
accessibile prima del pensionamento.

\subsection{La diversificazione dei rischi previdenziali:
            rischio demografico (I pilastro), rischio di mercato
            (II pilastro) e rischio inflazione (III pilastro)}
\label{subsec:diversificazione_rischi}

Ogni pilastro previdenziale è esposto a una classe specifica di rischio
sistemico che non è eliminabile internamente al pilastro stesso. La strategia
ottimale per il lavoratore NDC è quindi costruire un portafoglio previdenziale
\emph{diversificato per tipologia di rischio}, non per tipologia di strumento.

\begin{table}[htbp]
\centering
\caption{Matrice dei rischi sistemici per pilastro previdenziale}
\label{tab:rischi_pilastri}
\begin{tabular}{p{3.2cm}p{3.5cm}p{4.0cm}p{3.5cm}}
\toprule
 & \textbf{I Pilastro} (PAYG/NDC) &
   \textbf{II Pilastro} (Previdenza complementare) &
   \textbf{III Pilastro} (Risparmio privato) \\
\midrule
\textbf{Rischio principale} &
  Demografico e politico &
  Di mercato e di longevità &
  Inflazione e mercato \\
\textbf{Meccanismo} &
  Riduzione del rapporto lavoratori/pensionati, modifiche legislative
  dei coefficienti &
  Crollo dei mercati finanziari, revisione dei rendimenti attesi &
  Erosione del potere d'acquisto, crisi di liquidità \\
\textbf{Orizzonte del rischio} &
  Strutturale (30--50 anni) &
  Ciclico (7--15 anni) &
  Congiunturale (1--5 anni) \\
\textbf{Protezione offerta agli altri pilastri} &
  Nessuna (tutti i rischi rimangono al lavoratore) &
  Parziale (disaccoppiamento da g) &
  Alta liquidità in caso di crisi \\
\bottomrule
\end{tabular}
\smallskip\\
\footnotesize
Fonti: elaborazione propria su \cite{WorldBank1994, WorldBank2005, Bosi2018}.
\end{table}

La Tabella~\ref{tab:rischi_pilastri} mostra che i rischi dei tre pilastri
sono tra loro in larga misura \emph{non correlati}: una recessione demografica
che abbatte le pensioni NDC (rischio del I pilastro) non implica
necessariamente il crollo dei mercati finanziari; un'ondata di inflazione
che erode il risparmio privato (rischio del III pilastro) non modifica
i coefficienti di trasformazione del II pilastro. Diversificare tra i tre
pilastri non è, quindi, una metafora finanziaria applicata alla previdenza:
è la reazione razionale all'esistenza di rischi sistemici indipendenti che
colpiscono classi diverse di asset previdenziali~\cite{WorldBank2005}.

\subsection{Il framework World Bank dei tre pilastri (1994)
            e il suo aggiornamento a cinque pilastri (2005):
            una visione integrata del patrimonio previdenziale}
\label{subsec:world_bank}

Il framework multi-pilastro fu sistematizzato dalla Banca Mondiale nel
rapporto \emph{Averting the Old Age Crisis} (1994), che proponeva un modello
a tre pilastri: un primo pilastro pubblico obbligatorio a ripartizione,
funzionale a garantire un reddito minimo di sussistenza; un secondo pilastro
obbligatorio a capitalizzazione individuale, finalizzato al mantenimento del
tenore di vita; e un terzo pilastro volontario per i risparmi
aggiuntivi~\cite{WorldBank1994}. Questo framework, adottato come
riferimento nelle riforme pensionistiche degli anni '90 in Europa orientale,
America Latina e Africa, aveva il pregio della chiarezza ma il difetto di
ignorare i rischi di mercato del secondo pilastro obbligatorio e la
protezione sociale dei più vulnerabili.

L'aggiornamento del 2005 -- \emph{Old Age Income Support in the 21st Century}
-- ha esteso il modello a cinque pilastri~\cite{WorldBank2005}:
%
\begin{enumerate}
  \item[\textbf{Pilastro 0}] -- Pensione di cittadinanza non contributiva
    (safety net universale, finanziata dalla fiscalità generale): protegge
    chi non ha mai lavorato o ha carriere troppo brevi per maturare diritti
    contributivi.
  \item[\textbf{Pilastro 1}] -- Sistema pubblico obbligatorio a ripartizione
    (NDC o DB): l'attuale primo pilastro italiano.
  \item[\textbf{Pilastro 2}] -- Sistema complementare obbligatorio o
    quasi-obbligatorio a capitalizzazione: il secondo pilastro, in Italia
    ancora volontario.
  \item[\textbf{Pilastro 3}] -- Risparmio privato volontario (ETF, PAC,
    polizze vita, immobili): il terzo pilastro informale discusso nelle
    sezioni precedenti.
  \item[\textbf{Pilastro 4}] -- Support non finanziario: reti familiari,
    supporto informale, accesso a cure sanitarie, housing. Formalmente fuori
    dal perimetro previdenziale ma cruciale per il benessere nella vecchiaia.
\end{enumerate}
%
Il sistema italiano del 2025 copre in modo adeguato solo il Pilastro~1 e,
parzialmente, il Pilastro~0 tramite l'assegno sociale. Il Pilastro~2 è
strutturalmente sottosviluppato (adesione al 38\%), il Pilastro~3 è
lasciato alla discrezionalità individuale senza alcun supporto informativo
pubblico, e il Pilastro~4 è invisibile alle politiche
pubbliche~\cite{WorldBank2005, MEFRGS2025}. Un sistema che si regge quasi
esclusivamente sul Pilastro~1 in un contesto di transizione demografica
e bassa crescita non è, in senso tecnico, un sistema previdenziale: è
un trasferimento fiscale intergenerazionale che scarica l'intero rischio
sulle generazioni future.

\subsection{Lo scenario di povertà relativa in età anziana:
            il profilo del lavoratore autonomo/precario Gen Z
            che si affida solo al primo pilastro}
\label{subsec:scenario_poverta}

Per rendere concreta la portata del problema, si considera il profilo di
Marco nella sua configurazione più vulnerabile: lavoratore autonomo puro
(aliquota contributiva \(c = 0{,}24\)), carriera con tre anni di interruzione
contributiva (di cui uno per formazione, uno per disoccupazione frizionale
e uno per cura familiare non retribuita), ingresso nel mercato del lavoro
a 29 anni, pensionamento a 67 anni (2060), speranza di vita residua di
21 anni. Scenario macroeconomico pessimistico (\(g = 1{,}5\%\)).

\begin{table}[htbp]
\centering
\caption{Simulazione di scenario pessimistico: Marco, lavoratore autonomo
         con carriera discontinua, solo primo pilastro
         (\(c=0{,}24\), \(S_0=\text{\euro}\,20.000\), \(m=1{,}5\%\),
          \(g=1{,}5\%\), \(L=35\) anni effettivi, \(e(67)=21\) anni)}
\label{tab:scenario_poverta}
\begin{tabular}{lrl}
\toprule
Variabile & Valore & Note \\
\midrule
Anni di contribuzione effettivi & 35 & 38 attesi $-$ 3 anni gap \\
Salario iniziale & \euro\,20.000 & p.IVA regime forfettario \\
Salario finale stimato & \euro\,29.800 & con \(m=1{,}5\%\) per 35 anni \\
Montante contributivo & \euro\,320.400 & formula~\eqref{eq:mc_m} \\
Coefficiente trasformazione & 5,723\% & valore 2023--2024 \\
Pensione annua lorda & \euro\,18.344 & \euro\,1.529/mese \\
Tasso di sostituzione lordo & $\approx 47\%$ & vs last salary \euro\,29.800 \\
Tasso di sostituzione netto & $\approx 52\%$ & su reddito netto \\
\midrule
Soglia povertà relativa (ISTAT 2024) & \euro\,13.428 & persona sola\\
Gap vs soglia povertà relativa & +\euro\,4.916 & margine esiguo \\
Soglia mantenimento tenore di vita (70\%) & \euro\,20.860 & \\
Gap vs soglia 70\% & $-$\euro\,2.516 & deficit strutturale \\
\bottomrule
\end{tabular}
\smallskip\\
\footnotesize
Elaborazione propria. Fonti: \cite{MEFRGS2025, ISTAT2024, Bosi2018}.
\end{table}

La Tabella~\ref{tab:scenario_poverta} rivela una realtà scomoda: un
lavoratore autonomo della Gen Z con una carriera non eccezionale ma nemmeno
particolarmente penalizzata si avvicina alla pensione con \euro\,1.529
mensili lordi -- un importo che si colloca a \euro\,4.900 annui sopra la
soglia di povertà relativa ISTAT per una persona sola (\euro\,13.428 nel
2024), ma strutturalmente al di sotto della soglia del 70\% necessaria a
mantenere il tenore di vita~\cite{ISTAT2024, Bosi2018}. Non si tratta di
uno scenario catastrofico: si tratta del profilo atteso -- non di un caso
limite -- per milioni di lavoratori autonomi e precari della Generazione Z
italiana che non attivano il secondo e il terzo pilastro. Questo è il costo
reale, misurato in euro mensili, dell'inerzia previdenziale.

% ============================================================
\section{Educazione finanziaria previdenziale come bene pubblico}
\label{sec:educazione_finanziaria}
% ============================================================

Le sezioni precedenti hanno mostrato che esistono strategie razionalmente
superiori alla semplice dipendenza dal primo pilastro: l'investimento in
capitale umano, la diversificazione tra pilastri, la scelta dello strumento
complementare con il rapporto costo-rendimento più favorevole. La domanda
che occorre porsi è: perché queste strategie non vengono adottate dalla
maggioranza dei lavoratori italiani? La risposta standard dell'economia
neoclassica -- irrazionalità individuale -- è empiricamente inadeguata e
politicamente pericolosa, perché sposta la colpa sul singolo e solleva
lo Stato dalla propria responsabilità. La risposta corretta è più precisa:
il mercato dell'informazione previdenziale \emph{fallisce}, e questo
fallimento non è casuale ma strutturale.

\subsection{Il fallimento del mercato informativo previdenziale:
            perché la bassa adesione al secondo pilastro
            non è irrazionalità ma ignoranza strutturale}
\label{subsec:fallimento_informativo}

Tre caratteristiche strutturali del problema previdenziale creano le
condizioni per un fallimento sistematico del mercato dell'informazione.

\textbf{Prima caratteristica: l'orizzonte temporale estremo.} Il feedback
della scelta previdenziale arriva 40 anni dopo la scelta stessa. La
psicologia economica ha documentato sin dagli studi di Kahneman e Tversky
che gli individui sono sistematicamente miopi rispetto a conseguenze lontane
nel tempo (\emph{hyperbolic discounting}): attribuiscono un peso enorme alle
conseguenze immediate e quasi zero a quelle distanti~\cite{ThalerSunstein2008}.
Questa miopia temporale non è un difetto di carattere: è una caratteristica
cognitiva documentata neuroscientificamente, che non può essere corretta da
campagne di comunicazione generiche.

\textbf{Seconda caratteristica: l'opacità tecnica.} Il calcolo della pensione
NDC richiede la comprensione della formula~\eqref{eq:mc_m}, dei coefficienti
di trasformazione, dei tassi di capitalizzazione nozionale e della loro
dipendenza dalla crescita del PIL nominale. Secondo i dati dell'indagine
OCSE/INFE sulla literacy finanziaria (2023), solo il 32\% dei giovani
italiani sa cos'è un fondo pensione negoziale e solo il 18\% è in grado
di stimare approssimativamente la propria pensione
attesa~\cite{OECD2023, INPS2024}. Questo non è analfabetismo finanziario:
è il risultato atteso di un sistema che non ha mai investito nella
trasmissione delle informazioni previdenziali in forma comprensibile.

\textbf{Terza caratteristica: l'asimmetria informativa a favore dei
  produttori.} I fornitori di prodotti previdenziali (compagnie assicurative,
fondi, banche) hanno un incentivo commerciale a enfatizzare i vantaggi dei
propri prodotti e a occultare i costi. Le commissioni dei PIP, pari in media
al 2,17\% annuo (ISC), vengono presentate come percentuali di rendimento
rinunciate -- non come capitale perduto su quarant'anni~\cite{COVIP2024}.
In assenza di un'informazione pubblica neutrale e accessibile, il lavoratore
inesperto opera in condizioni di ignoranza strutturale: sceglie male non
perché sia irrazionale, ma perché le informazioni necessarie per scegliere
bene non gli vengono fornite~\cite{Bosi2018, ThalerSunstein2008}.

Solo un italiano su tre ha destinato il TFR al fondo pensione~\cite{COVIP2024,
INPS2024}. Questo dato non misura la preferenza dei lavoratori per il risparmio
forzato aziendale: misura l'esito di un fallimento informativo che lo Stato ha
la responsabilità e la capacità di correggere, a un costo marginale
quasi nullo.

\subsection{La busta arancione italiana (INPS ``La Mia Pensione''):
            perché non funziona e come renderla
            uno strumento di consapevolezza efficace}
\label{subsec:busta_arancione}

Il modello svedese della \emph{orange kuvert} (introdotto nel 1999) ha
dimostrato empiricamente che la comunicazione proattiva e periodica della
proiezione pensionistica personalizzata aumenta in modo significativo la
consapevolezza previdenziale e l'adesione al secondo pilastro: a cinque
anni dall'introduzione della busta arancione, la quota di lavoratori svedesi
iscritti al fondo pensione complementare è cresciuta di oltre 30 punti
percentuali~\cite{Sundén2006}. L'Italia ha replicato formalmente questo
strumento nel 2015 con la piattaforma INPS ``La Mia Pensione'', disponibile
sul portale \texttt{inps.it}. La replica, tuttavia, si è fermata alla forma
senza replicarne la sostanza.

I quattro difetti strutturali della piattaforma INPS possono essere
identificati con precisione:
%
\begin{enumerate}
  \item \textbf{Passività}: la proiezione esiste, ma non viene inviata
    proattivamente al lavoratore. Chi non accede spontaneamente al portale
    INPS -- la grande maggioranza -- non riceve mai la propria stima.
    Il tasso di utilizzo stimato è inferiore al 15\% dei lavoratori
    attivi~\cite{INPS2024}.
  \item \textbf{Non periodicità}: non esiste un invio annuale sistematico.
    L'informazione è disponibile \emph{on demand}, non per default.
  \item \textbf{Opacità delle proiezioni}: le stime sono presentate in
    termini nominali futuri, senza deflazionare per l'inflazione attesa
    e senza scenari alternativi (con/senza fondo pensione). Il lavoratore
    vede un numero che non sa interpretare.
  \item \textbf{Mancanza di comparazione}: nessuna simulazione mostra
    l'impatto dell'adesione al secondo pilastro sulla pensione attesa,
    che è l'informazione più direttamente utile per prendere
    una decisione~\cite{INPS2024, Sundén2006}.
\end{enumerate}

La proposta di riforma è tecnicamente banale e politicamente a basso costo.
Ogni anno, contestualmente all'elaborazione della Certificazione Unica,
l'INPS invia a ciascun lavoratore attivo -- per via digitale (MyINPS) o
cartacea per chi non ha domicilio digitale -- un estratto previdenziale di
due pagine che contenga: (i) il montante contributivo accumulato a fine
anno, (ii) la pensione attesa in tre scenari (pessimistico, base,
ottimistico), (iii) la stessa pensione attesa con e senza adesione a un
fondo pensione di categoria, (iv) la data stimata di maturazione dei
requisiti per la pensione di vecchiaia. Il costo marginale per lo Stato
è quasi zero (i dati esistono già nel sistema INPS); l'impatto atteso,
sulla base delle evidenze svedesi, è un aumento significativo dell'adesione
al secondo pilastro~\cite{Sundén2006, ThalerSunstein2008}.

\subsection{Le fintech previdenziali per la Gen Z:
            agenti di mercato che documentano un fallimento
            pubblico e aprono spazio a partnership pubblico-privato}
\label{subsec:fintech}

Nell'assenza di informazione pubblica adeguata, il mercato ha risposto
con soluzioni private. Le fintech europee specializzate in educazione
finanziaria e pianificazione previdenziale per i giovani lavoratori -- da
piattaforme di robo-advisory previdenziale a app di simulazione
pensionistica -- stanno colmando un vuoto informativo che lo Stato non
riesce a coprire. Questo fenomeno è un \emph{segnale di mercato}, non una
soluzione di mercato: documenta l'esistenza di una domanda insoddisfatta di
informazione previdenziale trasparente, accessibile e personalizzata da parte
di una generazione -- la Gen Z -- che non si fida del sistema pubblico e
non ha i mezzi per affidarsi a un consulente finanziario tradizionale.

Le fintech previdenziali non sono, e non possono essere, un sostituto della
policy pubblica per almeno due ragioni strutturali. Prima: raggiungono
prevalentemente lavoratori già finanziariamente consapevoli, lasciando
esclusi esattamente quelli più esposti al rischio di inadeguatezza
(lavoratori a basso reddito, discontinui, a bassa istruzione). Seconda:
operano su un modello commerciale che ne allinea gli incentivi al
trattenimento dell'utente sulla piattaforma, non necessariamente alla qualità
delle sue scelte previdenziali~\cite{ThalerSunstein2008}.

Ciò che le fintech dimostrano è che il modello UK del \emph{Money and
Pensions Service} -- agenzia pubblica indipendente che eroga consulenza
finanziaria e previdenziale gratuita in partnership con operatori privati
accreditati -- è replicabile anche in Italia a costi contenuti. Una
partnership pubblico-privato che affidi alla fintech accreditata la
front-end dell'informazione (app, simulatori, interfacce user-friendly) e
mantenga allo Stato la garanzia della neutralità e della completezza
delle informazioni potrebbe aumentare significativamente l'efficacia
dell'educazione finanziaria previdenziale, con costi inferiori a quelli
di qualsiasi campagna informativa pubblica tradizionale~\cite{WorldBank2005}.

\subsection{Il nudge previdenziale: auto-enrollment,
            contribuzione crescente e semplificazione della scelta
            come policy a costo quasi zero}
\label{subsec:nudge}

La proposta più elegante e meno costosa per correggere il fallimento
informativo viene dall'economia comportamentale. Richard Thaler e Cass
Sunstein, in \emph{Nudge: Improving Decisions About Health, Wealth, and
Happiness} (2008), hanno formalizzato il concetto di \emph{architettura
delle scelte}: la struttura in cui viene presentata una decisione influenza
sistematicamente l'esito della decisione stessa, indipendentemente dalle
preferenze dell'individuo~\cite{ThalerSunstein2008}. Il \emph{nudge}
previdenziale sfrutta questo principio per aumentare l'adesione al secondo
pilastro senza ricorrere alla coercizione.

Tre nudge sono direttamente applicabili al sistema pensionistico italiano,
con costi quasi nulli per lo Stato e impatti documentati da evidenze
internazionali robuste.

\paragraph{Nudge 1 -- Auto-enrollment (opt-out vs opt-in).}
Il sistema attuale prevede che il lavoratore si iscriva attivamente al
fondo pensione (sistema opt-in). Il silenzio-assenso del D.Lgs. 252/2005
riguarda solo il conferimento del TFR, non l'iscrizione al fondo. Un
sistema di auto-enrollment strutturale -- in cui il lavoratore viene
iscritto automaticamente al fondo negoziale di categoria al momento
dell'assunzione, con possibilità di recesso (opt-out) entro 30 giorni --
sfrutta lo \emph{status quo bias}: la tendenza a mantenere l'opzione
di default senza prendere un'azione attiva per cambiarla~\cite{ThalerSunstein2008}.

L'evidenza del programma NEST nel Regno Unito (introdotto nel 2012) è
inequivocabile: nei tre anni successivi all'introduzione dell'auto-enrollment
obbligatorio per i datori con più di 250 dipendenti, il tasso di adesione
al secondo pilastro è passato dal 55\% al 90\%~\cite{DWP2017}.
La quasi totalità dell'aumento non è il risultato di nuove iscrizioni
deliberate: è il risultato dell'inerzia -- gli stessi individui che non
si erano mai iscritti attivamente non si sono nemmeno disiscritti quando
l'iscrizione è diventata il default.

\paragraph{Nudge 2 -- Save More Tomorrow (contribuzione crescente).}
Thaler e Benartzi (2004) hanno proposto e sperimentato il programma
\emph{Save More Tomorrow} (SMarT): i lavoratori si impegnano in anticipo
ad aumentare la propria contribuzione al fondo pensione in coincidenza con
i futuri aumenti di stipendio~\cite{ThalerBenartzi2004}. Il meccanismo
sfrutta due bias comportamentali: la difficoltà di ridurre il reddito
corrente (se il contributo aumenta solo quando lo stipendio aumenta, il
netto in busta paga non scende mai) e la minore resistenza agli impegni
futuri rispetto a quelli immediati. In quattro cicli di applicazione nelle
aziende americane coinvolte, il tasso medio di contribuzione al fondo
pensione è passato dal 3,5\% al 13,6\% della retribuzione in 28
mesi~\cite{ThalerBenartzi2004}.

\paragraph{Nudge 3 -- Semplificazione della scelta (meno opzioni,
  scelta migliore).}
Barry Schwartz, nel \emph{Paradox of Choice} (2004), ha documentato che
un eccesso di opzioni disponibili non aumenta la qualità delle scelte
ma la paralisi decisionale~\cite{Schwartz2004}. Il mercato italiano della
previdenza complementare offre attualmente oltre 300 fondi tra negoziali,
aperti e PIP, con strutture di costo e di rendimento difficilmente
comparabili da un non-esperto. L'introduzione di un \emph{fondo di
default} per categoria -- un fondo a gestione passiva a basso costo
assegnato automaticamente a chi non sceglie attivamente -- ridurrebbe
la paralisi decisionale e orienterebbe la massa degli iscritti verso
il veicolo con il miglior rapporto costo-rendimento, senza limitare
la libertà di chi vuole scegliere altrimenti.

\begin{table}[htbp]
\centering
\caption{Tre nudge previdenziali applicabili al sistema italiano:
         meccanismo, evidenza internazionale e costo per lo Stato}
\label{tab:nudge}
\begin{tabular}{p{3.5cm}p{4.0cm}p{3.5cm}p{2.5cm}}
\toprule
\textbf{Nudge} & \textbf{Meccanismo comportamentale} &
  \textbf{Evidenza internazionale} & \textbf{Costo per lo Stato} \\
\midrule
Auto-enrollment &
  Status quo bias: l'opt-out è più raro dell'opt-in &
  UK NEST: adesione 55\%$\to$90\% in 3 anni~\cite{DWP2017} &
  Quasi zero (modifica normativa) \\
Save More Tomorrow &
  Resistenza alla perdita immediata mitigata da impegno futuro &
  USA: contribuzione 3,5\%$\to$13,6\% in 28 mesi~\cite{ThalerBenartzi2004} &
  Zero (accordi collettivi) \\
Fondo di default &
  Paradosso della scelta: meno opzioni, più adesione &
  Svezia (2000): 90\% dei nuovi iscritti non cambia il default~\cite{Sundén2006} &
  Quasi zero (modifica regolamentare) \\
\bottomrule
\end{tabular}
\smallskip\\
\footnotesize
Fonti: \cite{ThalerSunstein2008, ThalerBenartzi2004, DWP2017, Sundén2006}.
\end{table}

I tre nudge della Tabella~\ref{tab:nudge} condividono una proprietà
fondamentale: non richiedono né coercizione né spesa pubblica aggiuntiva.
Richiedono soltanto la volontà politica di riformare l'architettura delle
scelte previdenziali in un sistema che, allo stato attuale, è progettato
come se i lavoratori fossero agenti perfettamente razionali, perfettamente
informati e immuni dalla miopia temporale. La letteratura comportamentale
ha dimostrato con rigore che nessun lavoratore -- a prescindere dal livello
di istruzione -- è immune da questi bias. Ignorarlo non è rigore liberale:
è negligenza istituzionale~\cite{ThalerSunstein2008, Bosi2018}.